\section{Introduction}
\label{sec:introduction}

Mobile dual-arm robots are increasingly expected to service multiple workspaces in laboratories, restaurants, warehouses, and human-assistance environments. Recent studies on mobile manipulation and task-and-motion planning show that such robots must coordinate navigation, base placement, arm feasibility, and task execution rather than solve them as isolated modules \cite{sandakalum2022motion,zhao2024survey,zhang2022grop,gu2023multiskill,fu2024mobilealoha}. In multi-table service settings, the robot must navigate between tables or benches while executing manipulation tasks distributed over tabletop surfaces. The central planning question is therefore not only how to find a collision-free base path, but also how to choose base poses from which the arms can actually reach the assigned objects or task points.

The difficulty arises because tabletop task coordinates are not valid navigation goals for the mobile base. A task point may be inside the arm workspace from one side of a table but unreachable from another side, while the mobile base must remain outside the table footprint and satisfy clearance and orientation constraints. A planner that directly connects task locations in the ground plane can report a short path even when no feasible arm motion exists after arrival. This mismatch is especially severe for dual-arm platforms, where left-arm and right-arm reachability, table geometry, and approach direction jointly determine whether a docking pose is useful.

Existing motion-planning methods solve important parts of the problem but do not fully address this decision layer. Grid and graph search methods such as Dijkstra's algorithm and A* provide reliable global path planning in known maps \cite{dijkstra1959note,hart1968formal,liu2023pathplanning}, and any-angle or symmetry-reduced variants can reduce path length or search effort \cite{nash2007theta,harabor2011online}. Mobile-manipulation studies have examined robot placement, reachability priors, approach uncertainty, and base-pose sequence planning \cite{reister2022combining,jauhri2022reachability,yamazaki2022approaching,naik2024basenet}. However, tightly coupled planning over all base, arm, and task-order variables remains costly for multi-table scenes. Coverage and viewpoint-selection methods reason about covering spatial targets, yet they often lack explicit dual-arm reachability constraints and route-level integration.

The first limitation is that path optimality can conflict with manipulation feasibility. A nearest docking or shortest-route choice may reduce travel distance while leaving many tabletop tasks uncovered. The experiments in this paper show this failure mode directly: nearest docking achieves only 50.00\% completion in the main scenario and drops to 0.00\% in hard mixed task distributions. Thus, for mobile manipulation, coverage and reachability should be treated as feasibility conditions before route length is minimized.

The second limitation is that a single docking pose is often insufficient for a complex table. Tasks can lie on different sides or near reachability boundaries, and no single base pose may cover all of them. Enforcing one docking pose per table converts a physically feasible multi-stop operation into an infeasible planning problem. A planner therefore needs a mechanism that allows multiple docking poses to jointly cover the same tabletop region.

The third limitation is that global path planning and docking selection are often evaluated separately. A mobile-manipulation planner should preserve a clean interface between local manipulation feasibility and global navigation so that different global planners can be benchmarked under the same selected docking set. This modularity is important for reproducibility, algorithm comparison, and deployment in different map or robot settings.

Based on these observations, this paper proposes RCH-MT-MMP, a reachability-constrained hierarchical planning framework that transforms tabletop tasks into reachable docking-pose coverage decisions and then connects the selected docking poses through a constrained global routing layer.

The contributions are as follows.

\begin{enumerate}
    \item A constrained hierarchical formulation is introduced for multi-table dual-arm mobile manipulation. The formulation separates local docking feasibility from global route connection while keeping them in a single objective that penalizes uncovered tasks, docking cost, and route length.
    \item A candidate docking-pose evaluation method is developed using task coverage, dual-arm reachability, safety distance, orientation consistency, and approach cost. This scoring rule converts geometric task distributions into ranked service poses suitable for dual-arm operation.
    \item A coverage-aware multi-docking selection strategy is designed so that one table may be covered by several docking poses. Under a mild candidate-coverage condition, the greedy selection loop terminates with full coverage in finite time.
    \item A unified experimental pipeline is provided for main scenes, hard task distributions, ablations, random-scale studies, noise and stress tests, small-scale oracle references, and global-planner benchmarks. In submission-level experiments, the full method maintains 100.00\% completion, while single-docking and nearest-docking baselines leave substantial tasks uncovered.
\end{enumerate}

The rest of the paper is organized as follows. Section~\ref{sec:related} reviews related work. Section~\ref{sec:problem} formulates the constrained docking and routing problem. Section~\ref{sec:method} presents the proposed framework. Section~\ref{sec:analysis} gives feasibility and complexity analysis. Section~\ref{sec:experiments} reports the experimental evidence, and Section~\ref{sec:conclusion} concludes the paper.
